% Member 2 - final report, II.B Style Transfer and AdaIN
\subsection{Style Transfer and AdaIN}
\label{sec:adain}

The appearance change we need is a large one: it must move an image from the
photographic domain to the painterly one while keeping the object
recognisable. Two families of tools can do this. Image-to-image translation
methods such as CycleGAN~\\cite{zhu2017unpaired} learn a mapping between two image
collections from unpaired examples. They can produce striking changes, but
they need a target collection to learn from and they change the
\emph{content} as well as the appearance, so the class label of the training
image may no longer be justified. Neural style transfer
\\cite{gatys2016style} and AdaIN~\\cite{huang2017adain} instead combine the
\emph{content} of one image with the \emph{style} of another, which keeps the
content exactly fixed and is therefore well suited to label-preserving data
augmentation.

AdaIN is the instance of that family we use. It aligns the channel-wise
mean and variance of the content feature map with those of the style feature
map:

\begin{equation}
\mathrm{AdaIN}(F_c, F_s) = \sigma(F_s)\left(\frac{F_c - \mu(F_c)}{\sigma(F_c)}\right) + \mu(F_s),
\label{eq:adain}
\end{equation}

where $\mu$ and $\sigma$ are taken per channel over the spatial dimensions.
The encoder--decoder network is the pretrained one introduced with AdaIN: a
VGG-19 encoder trained for perceptual loss and a mirrored decoder, used
frozen, with no training of the generator at any point. The normalisation
itself, and the feed-forward stylisation pipeline that AdaIN belongs to,
were developed in~\cite{ulyanov2016instance,ulyanov2017improved}; AdaIN
replaced the fixed scale and shift of instance normalisation by the
statistics of the style image, which is what makes the strength of the
result interpretable. We use this
pretrained network rather than training our own because it keeps the
comparison honest: the only difference between our method and the AdaIN
baseline is the rule by which the two feature maps are combined, not the
quality of the generator.

Equation~\eqref{eq:adain} mixes the two feature maps with a single constant
weight, so the stylisation strength is uniform over the whole image. That
uniformity is exactly what our method relaxes: we replace the constant with
a per-location weight, so that locations whose structure is worth protecting
receive less style. Section~\ref{sec:structure} describes that
modulation, and Section~\ref{sec:strength_relation} shows empirically what it does
and does not buy.

Two properties of AdaIN make it a fair and useful baseline for this study.
First, it is deterministic and training-free, so every content--style pair
always yields the same image; the augmentation is therefore reproducible and
can be regenerated exactly, which is what makes the five-seed comparison
below meaningful. Second, the amount of style it applies is controlled by a
single interpretable quantity, the strength $\alpha$ with which we blend
between the original and the AdaIN feature map. That makes AdaIN the
natural control for our method: it applies the same operation, with the same
network, at a strength we can match exactly.
